\section{Metodología de Evaluación}

Para evaluar con precisión los resultados y el impacto de la implementación de Odoo en Industria Textil Atlas E.I.R.L., se aplicó un enfoque de análisis integral orientado a medir la transición desde los procesos logísticos y comerciales manuales hacia el ecosistema digital integrado. La evaluación se estructuró bajo los siguientes cuatro pilares:

\begin{itemize}
    \item \textbf{Comparación de procesos (Antes vs. Después):} Se contrastó el flujo de trabajo previo ---basado en el registro manual en cuadernos de venta de mostrador, búsqueda física improvisada de productos en tienda y pedidos reactivos de compra a confeccionistas--- con el nuevo flujo unificado en Odoo. Esto permitió validar la consistencia en el control del almacén físico.
    \item \textbf{Análisis de productividad y eficiencia comercial:} Se midieron los tiempos de atención directa en el Punto de Venta (POS) de San Camilo y el tiempo de respuesta comercial para la elaboración de presupuestos mayoristas en el módulo de Ventas, comprobando la rapidez y exactitud del cálculo automático de importes.
    \item \textbf{Evaluación de Indicadores de Desempeño (KPIs):} Se confrontaron las métricas de la situación inicial extraídas del diagnóstico de madurez digital con los valores reales obtenidos durante la marcha blanca del sistema. Se analizaron variables como el tiempo promedio de despacho, la tasa de error por discrepancia de stock y el nivel de compras automatizadas.
    \item \textbf{Observación directa de adopción y capacitación:} Se evaluó la interacción diaria del personal de tienda y almacén con Odoo en base al nivel de uso del POS, registro de recepciones de stock y el control gerencial exclusivo de Víctor Huayllaro sobre costos y listas de precios, garantizando que no se mantengan procesos analógicos paralelos.
\end{itemize}
